\centerline{\bf \Huge Начало Чётности I}

\begin{flushright}
        \scriptsize{Чётность не меняют только чётные числа - этим они и хороши. Начальник}
\end{flushright}

\problem{\textbf{1}} Чётной или нечётной будет сумма 2025 чётных чисел? А 2025 нечётных?

\problem{\textbf{2}} IShowSpeed говорит, что знает 2024 числа, таких что их сумма и произведение нечётны. Прав ли он?

\problem{\textbf{3}} В ряд выписаны числа от 1 до 10. Можно ли между ними расставить знаки плюс и минус так, чтобы получился 0?

\problem{\textbf{4}} Сложили 15 натуральных чисел и получили 2025. Сколько среди них может быть нечётных? Перечислите все варианты!

\problem{\textbf{5}} Алина купила в магазине 20 тетрадей, 2 альбома для рисования, несколько карандашей по 6 р. 20 коп. и несколько ластиков по 4 рубля. Ей сказали, что в кассу следует уплатить 55 рублей 65 копеек. Алина попросила пересчитать стоимость покупки и ошибка была устранена. Как Алина догадалась, что она была допущена?

\problem{\textbf{5}} Конь вышел с поля a1 и через несколько ходов вернулся на него. Чётное или нечётное количество ходов он сделал?

\problem{\textbf{7}} 13 детей сели за круглый стол и договорились,
что мальчики будут врать девочкам, а друг другу говорить правду, а девочки, наоборот, будут
врать мальчикам, а друг другу говорить правду. Один из детей сказал своему правому соседу: «Большинство из нас мальчики». Тот сказал своему правому соседу: «Большинство из нас
девочки», а он своему соседу справа: «Большинство из нас мальчики», а тот своему: «Большинство из нас девочки» и так далее, пока последний ребёнок не сказал первому: «Большинство
из нас мальчики». Сколько мальчиков было за столом?